\documentclass[11pt]{article}
\usepackage[T1]{fontenc}
\usepackage[utf8]{inputenc}
\usepackage{lmodern,amsmath,amssymb,amsthm,mathtools}
\usepackage[margin=1in]{geometry}
\usepackage[hidelinks]{hyperref}
\newtheorem{theorem}{Theorem}
\newtheorem{lemma}[theorem]{Lemma}
\newtheorem{proposition}[theorem]{Proposition}
\newcommand{\R}{\mathbb R}\newcommand{\Z}{\mathbb Z}
\newcommand{\C}{\mathbb C}\newcommand{\N}{\mathbb N}
\newcommand{\ee}{\mathrm e}\newcommand{\Lie}{\operatorname{Lie}}
\newcommand{\id}{\mathrm{id}}\newcommand{\Ad}{\operatorname{Ad}}
\title{The Section 12 semidirect-product join}
\author{Author proposal: /root/astra\_v14\_baseline\_gaussian}
\date{7 November 2026}
\begin{document}\maketitle

\noindent\textbf{Scope.} This project proof supplies the final geometric
construction on published GTZ pp.~1313--1315 from the explicit petal quotient
and comparison interface of LP04. It is an author proposal, not independent
acceptance. The analytic frequency input, rank iteration and other inverse
theorem interfaces remain separate. We neither invoke the final inverse
theorem to prove an intermediate construction nor import the unpublished
erratum as a theorem. All quotients below are right lattice quotients.

\section{Exact input and output}
Fix the data of LP04: integers $m\geq1$, $1\leq R\leq m$, fixed finite
sets of core and retained petal letters with positive integer weights,
and the weighted free nilpotent quotient
\[
 G=\exp\mathfrak g,\qquad
 \mathfrak g=\mathfrak c\oplus\mathfrak p.
\]
The homogeneous relations discard all words of weight greater than $m$,
all weight-$m$ words of length greater than $R$, and all words containing
at least two petals. The last relation holds at every weight.
Thus $\mathfrak p$ is an abelian ideal. Write $P=\exp\mathfrak p$,
$G_j=G_{(j,0)}$, $P_j=P\cap G_j$, $G_0=G_1=G$, and $P_0=P_1=P$.
Let $\Gamma$ be the lattice generated by the exponentials of the letters,
and $\Lambda=P\cap\Gamma$. LP04's Hall-coordinate lemma supplies the
exact integer lattice coordinates and rationality of $G_j,P_j$.
Put $T=G_{(m,R)}$, and let $\eta:T\to\R$ be the descended continuous
character integral on $T\cap\Gamma$.

There are $\ell$ retained petal letters. For $t\in V=\R^\ell$, let
$M_t:P\to P$ scale the summand with unique petal letter $a$ by $t_a$.
We use the following proved LP04 identities, for all real parameters:
\begin{equation}\label{eq:actions}
 M_t(p)M_u(p)=M_{t+u}(p),\qquad
 M_t(gpg^{-1})=gM_t(p)g^{-1},\qquad M_t(P_j)\subseteq P_j.
\end{equation}
For integral $z$, $M_z(\Lambda)\subseteq\Lambda$. These are
endomorphisms; $M_0$ is not asserted to be invertible.

The family to be represented has the form
\begin{equation}\label{eq:inputfamily}
 b_h(n)=g_0(n)M_{\{\alpha h\}}(g_1(n)),\qquad
 \widetilde\chi_h(n)=\widetilde F(b_h(n)\Gamma).
\end{equation}
Here $g_0:\Z\to G$ and $g_1:\Z\to P$ are polynomial for the ordinary
weight filtrations; $\alpha\in\R^\ell$ is unrestricted. Braces denote
one fixed half-open fundamental cube of $V/\Z^\ell$, for example
$[-1/2,1/2)^\ell$. The factors used in GTZ are ordered products
$e_a^{\xi_a\binom ni}$ for core letters of weight $i$, and the commuting
product $\prod_{a\ \mathrm{petal}}e_a^{\beta_a\binom n{w(a)}}$ for $g_1$.
Core coefficients are independent of $h$. Identifying an arbitrary input
family with this specific factorization is an upstream obligation; it is
not inferred from the names of the factors.
For the actual coefficients in (12.3)--(12.4), the algebraic identification
is exact: $M_t(\exp(uX_a))=\exp(t_a uX_a)$, and $M_t$ is a
homomorphism on the abelian petal group. Hence applying
$M_{\{\alpha h\}}$ to the displayed $g_1$ gives precisely the product
in (12.4). No reordering past a core factor is involved. The hypotheses
that produce these coefficient data are still the separate analytic input.

The function $\widetilde F$ is a choice made below, before applying the
LP04 comparison. It is not an arbitrary previously fixed function having
only $T$-equivariance. This order of choices is necessary.

\begin{theorem}\label{thm:join}
One can choose $\widetilde F$ on the fixed $G/\Gamma$, with unit vector
norm and the required $T$-frequency $\eta$, such that the following holds.
There is a fixed connected simply connected nilpotent group $H$, a lattice
$\Delta$, and a rational taut $\N^2$-filtration supported on the downset
below $(1,m)$. The nilpotency class is at most $m+1$. There is a finite
list of fixed smooth unit-vector functions on $H/\Delta$, all with the
same top vertical character, with fixed finite Lipschitz bounds.

For every $\alpha,g_0,g_1$ as above and every finite nonempty parameter set
$S\subseteq\Z$, some $S'\subseteq S$ has $|S'|\geq |S|/J$, where the
fixed integer $J$ depends only on $\ell$, and one function $F'$ in this
list satisfies
\begin{equation}\label{eq:exactoutput}
 F'(O(h,n))=(\widetilde\chi_h(n),0)
 \quad(h\in S',\ n\in\Z).
\end{equation}
Here $O$ is a polynomial orbit of multidegree at most $(1,m)$.
The zero padding is a fixed isometric coordinate inclusion, so projection
onto the original coordinates gives literal equality with
$\widetilde\chi_h$. No scalar component is divided by.

All constants, manifolds and functions are independent of $\alpha$,
the polynomial coefficients, $S$ and $N$. The conclusion transfers to the
full internal ultrapower domains and preserves a dense set of limit
parameters. It supplies this construction and its LP04 comparison join;
it does not by itself establish the other premises of Theorem 7.1 or GI.
\end{theorem}

\section{Extend the character before choosing the test}
\begin{lemma}\label{lem:extension}
The character $\eta$ extends to a continuous character
$\bar\eta:G_m\to\R$ integral on $G_m\cap\Gamma$.
\end{lemma}
\begin{proof}
The weight-$m$ group is central. Its surviving Hall commutators have
weight $m$ and lengths at most $R$. At this weight the logarithm of each
group commutator is exactly its Lie bracket: every correction has greater
weight and vanishes. By LP04's exact Hall coordinates, these logarithms
form a real basis and their integer span is precisely the logarithm of
$G_m\cap\Gamma$. The subset of length $R$ is a basis of $T$, and its
integer span is $T\cap\Gamma$. Assign to those basis elements the values
of $\eta$, and assign zero to all remaining weight-$m$ basis elements.
These assigned values are integers. Linear extension on this central
vector group proves the assertion. The construction depends only on the
fixed Hall ordering and $\eta$.
\end{proof}

The central torus $G_m/(G_m\cap\Gamma)$ acts freely on $G/\Gamma$.
For any integral character of this torus there is a smooth equivariant
unit-vector function of finite dimension. Here is the construction we use
twice. Choose finitely many local smooth sections of the torus quotient
and real smooth functions $a_j$ supported in their domains with
$\sum a_j^2>0$ everywhere. Such functions can be obtained by smooth
Euclidean bumps in a finite collection of relatively compact local charts
whose smaller charts cover the compact base. Replace $a_j$ by
$a_j/(\sum a_i^2)^{1/2}$. In the $j$th local trivialization $(x,z)$ define
the $j$th component to be $a_j(x)\ee(\bar\eta(z))$ and extend it by zero.
The support lies inside its chart. Thus each component is globally smooth,
has the stated vertical transformation law, and the squared vector norm
is $\sum a_j^2=1$. Local sections can be seen directly: choose a linear
complement to the central Lie subspace. Centrality gives
$\exp(X+Z)=\exp X\exp Z$, so the complement gives smooth local lifts
from the group quotient. The image lattice is discrete by the Hall
coordinates with the central weight-$m$ coordinates deleted. Small
coordinate neighborhoods with disjoint lattice translates give the same
local lifts on the compact quotient.

Choose this function as $\widetilde F$. It has the original frequency
$\eta$ on $T$ and now also the required frequency on all of $G_m$.
Every cancellation in LP04's comparison uses only the restriction to $T$,
so this choice preserves that proof. The comparison does not require a
specific previously selected $\widetilde F$.

\paragraph{Adverse test.}
$T$ can be strictly smaller than $P_m$. For example, with $m=R=2$
and just one petal generator of weight two, $T$ is trivial whereas
$G_m=P\simeq\R$. The unit-vector function
$(\cos 2\pi x,\sin 2\pi x)$ on $\R/\Z$ satisfies the vacuous $T$-law
but has no scalar vertical character for all translations. Consequently
$T$-equivariance alone cannot be cited as the final multidegree frequency.
The integral extension and new choice above are the repair, not an
assertion that every old test already had the stronger property.

\section{Both semidirect products and their signs}
Set $K=G\ltimes P$ with multiplication
\begin{equation}\label{eq:Klaw}
 (g,p)(g',p')=(gg',\,(g')^{-1}pg'\,p').
\end{equation}
Associativity follows either directly from conjugation or from the
injective homomorphism $(g,p)\mapsto(g,gp)$ into $G\times G$.
Its image is the closed subgroup of pairs agreeing modulo $P$.
The inverse is $(g^{-1},gp^{-1}g^{-1})$. Its underlying manifold is
$G\times P$, so $K$ is connected and simply connected.

Define $\rho_t(g,p)=(gM_t(p),p)$. To check multiplication, applying
$\rho_t$ to \eqref{eq:Klaw} gives first component
$gM_t(p)g'M_t(p')$, by \eqref{eq:actions}. Multiplying the two images
has the same first component. Its second is
\[
 M_t(p')^{-1}(g')^{-1}pg'M_t(p')p'=(g')^{-1}pg'p',
\]
since all terms surrounding $p'$ are in the abelian group $P$.
Thus $\rho_t$ is a homomorphism. Further,
$\rho_t\rho_u=\rho_{t+u}$ and $\rho_0=\id$, so $\rho_t$ is an
automorphism with inverse $\rho_{-t}$. This conclusion uses the additive
shear action, not invertibility of $M_t$.

Set $H=V\ltimes_\rho K$ with the right-action convention
\begin{equation}\label{eq:Hlaw}
 (t,k)(t',k')=(t+t',\rho_{t'}(k)k').
\end{equation}
In three coordinates this reads
\[
 (t,g,p)(t',g',p')
 =(t+t',\,gM_{t'}(p)g',\,(g')^{-1}pg'p').
\]
The action identity proves associativity, and the inverse of $(t,k)$ is
$(-t,\rho_{-t}(k^{-1}))$. Smoothness follows from the polynomial BCH
formulas and the linear maps on $\mathfrak p$. The manifold
$V\times G\times P$ is connected and simply connected.

For a complete nilpotence and filtration check, use Lie coordinates
$(t,X,Y)\in V\oplus\mathfrak g\oplus\mathfrak p$. The bracket is
\begin{equation}\label{eq:bracket}
 [(t,X,Y),(t',X',Y')]
 =\bigl(0,\,[X,X']+M_{t'}Y-M_tY',\,[X,Y']+[Y,X']\bigr).
\end{equation}
Indeed the embedding of $K$ has derivative $(X,Y)\mapsto(X,X+Y)$,
giving the two bracket terms of $K$. The derivative of $\rho_t$ is
$(X,Y)\mapsto(X+M_tY,Y)$. Conjugation by $(t,1)$ on $K$ is
$\rho_{-t}$ in convention \eqref{eq:Hlaw}, giving the signs displayed.

Assign bidegree $(0,i)$ to core words of weight $i$ in the $X$ copy,
$(1,i)$ to petal words of weight $i$ in the $X$ copy, $(0,i)$ to petal
words of weight $i$ in the $Y$ copy, and $(1,0)$ to $V$.
Every nonzero term of \eqref{eq:bracket} has the sum of the input
bidegrees. For example, $[V,Y_i]$ lies in the first petal copy of degree
$(1,i)$, and brackets involving two petal words vanish. All weights are
at most $m$ and the first degree is at most one. Each generator has
positive total degree, so every bracket of length greater than $m+1$
vanishes. This proves nilpotence of $H$ of class at most $m+1$.

\section{The complete filtration and the lattice}
Define $H_{a,j}=1$ for every $a>1$ and every $j\geq0$.
The remaining positive-$j$ formulas, together with the zero-index groups, are
\begin{align*}
 H_{0,0}&=H, & H_{a,j}&=1\quad(a>1),\\
 H_{1,0}&=\{(t,p,1):t\in V,p\in P\},&
 H_{1,j}&=\{(0,p,1):p\in P_j\},\\
 H_{0,j}&=\{(0,g,p):g\in G_j,p\in P_j\}.
\end{align*}
For $j>m$ these groups are trivial. These definitions are also the sums
of the homogeneous Lie spaces whose bidegrees dominate the index.
Thus the groups are connected, descending and normal; brackets lie in
the sum-index subgroup by \eqref{eq:bracket}, and the BCH commutator
terms do also. To spell out the boundary cases, two $H_{1,0}$ elements
commute, $[H_{1,0},H_{0,j}]\subseteq H_{1,j}$, and
$[H_{1,i},H_{0,j}]\subseteq H_{1,i+j}$. Brackets with index zero are
covered by normality. The groups $H_{1,0}$ and $H_{0,1}=K$ generate
$H$, which proves tautness. In particular
\begin{equation}\label{eq:top}
 Z_*:=H_{1,m}=\{(0,p,1):p\in P_m\}
\end{equation}
is central. Importantly $H_{0,1}$ contains the full $G_1=G$ in its first
group coordinate, not merely $P$.

Let $\Lambda_K=\Gamma\ltimes\Lambda$ with law \eqref{eq:Klaw}.
Since $P$ is normal, $\Gamma$ normalizes $\Lambda$. Therefore this is
a subgroup, including inverses. For $z\in\Z^\ell$, both $\rho_z$
and $\rho_{-z}$ preserve $\Lambda_K$, by the integral LP04 action.
Consequently
\[
 \Delta=\Z^\ell\ltimes_\rho\Lambda_K
       =\{(z,\gamma,\lambda):z\in\Z^\ell,
                   \gamma\in\Gamma,\lambda\in\Lambda\}
\]
is a subgroup of $H$. It is discrete in the product manifold topology.

Here is explicit cocompactness. Right-multiply $(t,k)$ by $(-z,1)$,
where $z$ reduces $t$ into $[0,1]^\ell$. This gives
$(t-z,\rho_{-z}k)$. For an arbitrary $(g,p)\in K$, choose
$\gamma\in\Gamma$ with $g\gamma^{-1}$ in a fixed compact
representative set for $G/\Gamma$. Right multiplication by
$(\gamma^{-1},1)$ gives $(g\gamma^{-1},\gamma p\gamma^{-1})$.
Now reduce the second coordinate into a fixed compact representative
set for $P/\Lambda$ by right multiplication by $(1,\lambda^{-1})$.
The resulting compact product set represents every right coset of
$\Lambda_K$. Applying this after the $t$ reduction gives a compact
product set representing every $\Delta$ coset. No bound on the
intermediate conjugate or on $z$ is needed.

The same argument inside $G_j\ltimes P_j$, using the already proved
lattices $\Gamma\cap G_j$ and $\Lambda\cap P_j$, proves rationality
of $H_{0,j}$. The intersections for $H_{1,j}$ are exactly
$\Lambda\cap P_j$, and for $H_{1,0}$ exactly $\Z^\ell\times\Lambda$.
These are cocompact by LP04 and the product construction. This checks
all filtered lattice intersections, including zero and trivial groups.
The topology, filtration, lattice and smooth metric on the compact
quotient are fixed before any orbit coefficients are chosen.

\section{Polynomial orbit and exact reduction}
Define
\begin{equation}\label{eq:orbit}
 O(h,n)=(0,g_0(n),g_1(n))(\alpha h,1,1)\Delta.
\end{equation}
We verify polynomiality in the full group before taking the quotient.
The source filtration on $D=\Z^2$ is
$D_{0,0}=\Z^2$, $D_{1,0}=\Z\times\{0\}$,
$D_{0,1}=\{0\}\times\Z$, and $D_{a,j}=\{0\}$ at every other
nonzero index. It is an $\N^2$-filtered abelian group.
The embeddings $g\mapsto(0,g,1)$ and $p\mapsto(0,1,p)$ are
homomorphisms taking the weight-$j$ groups to $H_{0,j}$.
Each independent-of-$h$ factor is therefore polynomial on $\Z^2$ with
its coordinate multidegree filtration: $h$ differences vanish, and
$n$ differences have the asserted $H_{0,j}$ membership.
The map $(h,n)\mapsto(\alpha h,1,1)$ is a homomorphism in $h$ with
values in $H_{1,0}$, has zero second $h$ difference, and is constant
in $n$. Thus it too is polynomial for this same filtration.
Published Corollary B.4, GTZ p.~1334, says that polynomial maps between
the same index-filtered groups are closed under pointwise multiplication.
Here both source and target are the precise $\N^2$-filtered groups
just checked. Applying it to these three factors proves polynomiality
of \eqref{eq:orbit}, including all mixed derivatives and all integer inputs.
This is the only external polynomial-product theorem used in this join.

For the explicit inputs preceding Theorem~\ref{thm:join}, their univariate
polynomiality can also be checked before this use. An $r$-fold difference
of $e_a^{c\binom ni}$ is $e_a^{c\Delta^r\binom ni}$ because all its
values lie in one abelian one-parameter subgroup. It vanishes for $r>i$;
for $r\leq i$ it lies in $G_i\subseteq G_r$. This holds for arbitrary
integer increments, including negative ones. Products are polynomial by
the same Corollary B.4; for $g_1$ all factors already commute in $P$.
No Taylor expansion with a potentially unjustified centrality hypothesis
is imported for this check.

Write $z_h=\alpha h-\{\alpha h\}\in\Z^\ell$. Formula
\eqref{eq:Hlaw} gives a representative of \eqref{eq:orbit} as
\[
 (\alpha h,g_0(n)M_{\alpha h}(g_1(n)),g_1(n)).
\]
Right multiplication by $(-z_h,1,1)\in\Delta$ changes this to
\begin{equation}\label{eq:reduced}
 (\{\alpha h\},g_0(n)M_{\{\alpha h\}}(g_1(n)),g_1(n)).
\end{equation}
The equality follows from the additive identity in \eqref{eq:actions}.
It is exact for every $h,n$; there is no illegal commutation of a
fractional-part operation with a nonabelian product.

\section{Fixed charts, unit completion and uniform bounds}
Projection onto $t$ gives a smooth map
$\pi:H/\Delta\to V/\Z^\ell$. If a box $U\subset V$ has width less
than one in each coordinate, projection is injective on $U$. On the
corresponding open torus box there is an exact smooth trivialization
\begin{equation}\label{eq:chart}
 U\times K/\Lambda_K\longrightarrow\pi^{-1}(U\bmod\Z^\ell),
 \qquad(t,k\Lambda_K)\longmapsto(t,k)\Delta.
\end{equation}
Indeed two representatives in $U$ differ by an element of $\Delta$
only if their integer $t$ displacement is zero; their $k$ coordinates
then differ precisely by right multiplication in $\Lambda_K$.
Surjectivity follows by reducing the $t$ coordinate to its unique lift
in $U$. The formulas are smooth on this chart.

Choose finitely many fixed closed boxes $I_j$ covering the closed
fundamental cube $[-1/2,1/2]^\ell$, with fixed slightly larger boxes
$U_j$ of widths less than one. They may extend beyond the cube's
boundary; each $U_j$ still projects injectively. Choose a smooth
function $\theta_j$ on the torus which is zero on $I_j$ modulo integers
and $\pi/2$ outside a smaller open box with closure contained in $U_j$.
It can be constructed by a fixed Euclidean bump in $U_j$, equal to one
on $I_j$, followed by multiplication of its complement by $\pi/2$.
The chart margins are fixed positive real numbers, not functions of $N$.
At least one box contains $\{\alpha h\}$ for at least $|S|/J$ of
the parameters, where $J$ is the fixed number of boxes. Fix that box
and call this subfamily $S'$.

On its chart define
\[
 W(t,(g,p)\Lambda_K)=\widetilde F(g\Gamma).
\]
The projection $K/\Lambda_K\to G/\Gamma$ is well-defined and smooth,
since the first coordinate of \eqref{eq:Klaw} is a homomorphism and
maps $\Lambda_K$ into $\Gamma$. Thus $W$ is a smooth local unit
vector. Define a character on the actual group \eqref{eq:top} by
\[
 \eta'(0,p,1)=\bar\eta(p)\qquad(p\in P_m).
\]
It is integral on $Z_*\cap\Delta$ by Lemma~\ref{lem:extension}.
Left multiplication by $(0,p_*,1)\in Z_*$ changes the chart's first
group coordinate to $p_*g$, keeps $t$ fixed, and hence multiplies
$W$ by $\ee(\bar\eta(p_*))$. This uses equivariance on $P_m$,
not merely on $T\cap P$.

The torus $Z_* /(Z_*\cap\Delta)$ acts freely on $H/\Delta$:
centrality makes the stabilizer at every point exactly
$Z_*\cap\Delta$. It is compact by the rationality already proved.
The quotient by $Z_*$ has coordinates $V\times(G/P_m)\times P$,
and the image of $\Delta$ has coordinates
$\Z^\ell\times q(\Gamma)\times\Lambda$. Here $q(\Gamma)$ is
discrete by LP04's Hall-coordinate quotient, so the image lattice is
discrete and the quotient base is compact. Central complements and small
coordinate neighborhoods give local lifts exactly as in Section 2.
The finite local-section construction from that section therefore gives a
global smooth unit vector $W_0$ with frequency $\eta'$. Set
\begin{equation}\label{eq:completion}
 F'_j=(\cos(\theta_j\circ\pi)W,
                 \sin(\theta_j\circ\pi)W_0).
\end{equation}
The first component is extended by zero outside the chart. It vanishes
on a neighborhood of the chart boundary, so this extension is smooth.
Outside the chart the second coefficient is one. Inside, the squared
norm is $\cos^2\theta_j+\sin^2\theta_j=1$. Both components have the
same vertical frequency because $Z_*$ fixes $\pi$. Therefore $F'_j$
is a genuine unit-vector nilcharacter test, including outside the box.

For $h\in S'$, the representative \eqref{eq:reduced} uses the unique
lift $\{\alpha h\}\in I_j$. Its cutoff angle is zero, so
\eqref{eq:completion} gives \eqref{eq:exactoutput} for every integer $n$.
The second component is exactly zero there, not an approximation.

Every chart used has compact closure in a fixed larger chart, and its
fiber $K/\Lambda_K$ is compact. The displayed projections and their
derivatives are bounded on these sets. The finite bump functions,
$\widetilde F$, $W_0$, and the fixed metrics have bounded derivatives.
Taking the maximum over the finite list proves a finite Lipschitz bound
depending only on the fixed groups, character and charts. Unbounded
polynomial coefficients change only the orbit; they do not enter these
derivative bounds. This also proves fixed finite dimension. If $\ell=0$
the torus is a point, use one chart and $\theta=0$; all arguments still
apply. Trivial top groups or zero characters likewise cause no exception.
This completes the ordinary proof of Theorem~\ref{thm:join}.

\section{Internal domains and the LP04 consumer}
All groups, charts, test functions and characters above are standard
objects after fixing the selected dimension data and character. All
algebraic identities hold for every real parameter and every group point.
Taking representative sequences gives internal polynomial maps, including
for unbounded limit coefficients. The standard tests extend to the whole
ultrapower of the compact quotient. Their vertical laws then hold for
every limit element of $Z_*$ and every limit point, by transfer of the
actual universal identity. We are not inferring all-limit equivariance
from an identity checked only on standard points.

Finite pigeonhole over the fixed list of boxes preserves positive standard
density of a limit parameter family. Equivalently apply the ordinary
cardinality inequality to representatives and pass to the ultralimit.
After choosing the resulting fixed box, the nilcharacter
$\chi(h,n)=F'_j(O(h,n))$ is defined on all of $\!{}^*\Z^2$, not only
on the dense subfamily. Its exact restriction is the fixed isometric
inclusion of $\widetilde\chi_h$ for all limit integers $n$ there.

Choose $\widetilde F$ by Section 2 before invoking LP04's comparison
proposition. That proposition's top cancellation is unchanged because
$\bar\eta|_T=\eta$. It produces the same strictly lower degree--rank
comparison for $\chi_h\otimes\overline{\widetilde\chi_h}$.
For clarity, unit norm gives the exact component recovery identity
\[
 (\chi_h)_a=\sum_b (\widetilde\chi_h)_b
             \bigl((\chi_h)_a\overline{(\widetilde\chi_h)_b}\bigr).
\]
On the selected subfamily, replace $(\widetilde\chi_h)_b$ by the
corresponding component of the single $\chi(h,n)$. The terms in parentheses
are exactly the lower degree--rank sequences supplied by LP04. The sum
is finite and its coefficients are fixed; no division or nonvanishing of
an individual component is used. The zero padding in
\eqref{eq:exactoutput} changes none of these terms.

At $R=1$, lower degree--rank means degree at most $m-1$. At $R>1$
the subsequent rank iteration remains an upstream/downstream assembly
obligation; we do not incorrectly replace it by an immediate ordinary
degree drop. Likewise the frequency decomposition leading to
\eqref{eq:inputfamily}, the initial correlation, and the remaining
Theorem 7.1 induction premises are not established by this geometric join.

\section{Bounded related-use and consistency audit}
The search scope is the Section 12 construction and its immediate LP04
comparison and Theorem 7.1 consumers, published pp.~1308--1315. We also
check the polynomial-product supplier B.4 at p.~1334. This is an authored
bounded audit, not a claim of independent source-first coverage of GI.
\begin{enumerate}
\item The smaller degree--rank character domain versus the ordinary
top group is repaired by the integral extension and the new test choice.
The repeated final frequency on p.~1315 is explicitly restricted to
$P_m$ and uses $\bar\eta$. A favorable quotient-character review alone
does not supply this extension.
\item The illustrative $H_{0,1}$ on p.~1312 restricts its first group
coordinate to the petal subgroup, although the orbit contains the core
$e_2^{\beta n}$. For nonzero generic $\beta$ its $n$ difference is outside
that alleged group. The general p.~1314 definition uses full $G_j$;
we use and prove that definition, not the defective illustration.
\item Zero or negative scaling parameters do not make every $M_t$
invertible. The actual invertible automorphisms are the shears $\rho_t$,
whose inverse and integer lattice preservation are proved here.
\item The lattice in p.~1314 is a subgroup of $H$, not of the original
$G$. Every filtered lattice intersection and the exact right-coset
reduction are checked above; neither discreteness alone nor a real
endomorphism establishes the claim.
\item The cutoff on pp.~1312 and 1315 cannot alone be a unit-vector
test or guarantee equality on the selected box. Fixed inner boxes,
positive chart margins and \eqref{eq:completion} supply unit norm,
literal equality and coefficient-independent Lipschitz bounds.
\item In particular a variable box of arbitrarily small width would
not supply a fixed Lipschitz bound. The predetermined finite cover
removes this possible parameter dependence at a controlled density loss.
\item The p.~1308 claim that pigeonhole puts the fractional parts in a
short interval around zero is not valid for an arbitrary dense family.
For $\alpha=1/2$ and odd $h$, our fractional parts are all $-1/2$.
Our fixed cover allows a box at any location, including across the chosen
fundamental-cube boundary, and does not need a near-zero assertion.
\item Formula (12.11) is checked as right multiplication by an actual
lattice element; the displayed identities do not silently change the
object or introduce a different representative in the local test.
\item A lower rank comparison is not automatically a lower ordinary
degree comparison. The consumer paragraph above preserves the exact
LP04 conclusion and leaves the rank iteration to its actual proof.
\end{enumerate}
Earlier printed statements and earlier frozen project packets are retained;
this proposal does not relabel them as having contained the new argument.
Any other current copies of the smaller-domain frequency claim or the
illustrative filtration must be corrected or marked historical when this
proof is incorporated. The attached occurrence record identifies these
specific copies and the new canonical assertions for review.

\begin{thebibliography}{9}
\bibitem{GTZ} B.~Green, T.~Tao and T.~Ziegler,
\emph{An inverse theorem for the Gowers $U^{s+1}[N]$-norm},
Annals of Mathematics 176 (2012), 1231--1372,
Section 12 and Corollary B.4.
\url{https://doi.org/10.4007/annals.2012.176.2.11}.
\bibitem{LP04} Project proof by /root/astra\_v14\_assertions\_arithmetic,
\emph{GTZ-LP-04: the petal quotient and degree--rank construction},
lp04-001. Its exact manifest and proof hashes are recorded in the
accompanying source-use record. Its dependencies and independent-review
status remain separate.
\end{thebibliography}
\end{document}
